\oldpage[IV-2]{174}
La question étant locale sur $X$ et $Y$, on peut supposer que
$X=\operatorname{Spec}(A)$, $Y=\operatorname{Spec}(B)$, où $A$ et $B$ sont
deux anneaux intègres et intégralement clos, de corps des fractions $K$ et
$L$ respectivement; supposons d'abord que $A$ soit une $k$-algèbre de type
fini, donc $K$ une extension de type fini de $k$. Par platitude,
$A\otimes_k B$ est une sous-$k$-algèbre de $K\otimes_k L$; or
$K\otimes_k L$ est un anneau noethérien normal (en vertu de (6.7.4.1) ou
(6.14.2)), donc composé direct d'anneaux intègres $C_i$ en nombre fini, de
sorte que si $E_i$ est le corps des fractions de $C_i$, le composé direct
$E$ des $E_i$ est l'anneau total des fractions de $K\otimes_k L$; c'est
d'ailleurs aussi l'anneau total des fractions de $A\otimes_k B$, puisque
$K\otimes_k L$ est un anneau de fractions de $A\otimes_k B$. Cela étant,
on peut écrire
$A\otimes_k B=(A\otimes_k L)\cap(K\otimes_k B)$, et il résulte de (6.14.2)
que $A\otimes_k L$ et $K\otimes_k B$ sont normaux, donc $A\otimes_k B$ est
normal en vertu de (6.14.1.1).

Considérons maintenant le cas général; $A$ est alors réunion de la famille
filtrante croissante de ses sous-$k$-algèbres de type fini $A_\alpha$,
donc $A\otimes_k B$ est limite inductive des anneaux normaux
$A_\alpha\otimes_k B$. Pour prouver que $A\otimes_k B$ est normal, il
suffit donc (5.13.6) de prouver que toute composante irréductible de
$\operatorname{Spec}(A_\beta\otimes_k B)$ domine une composante
irréductible de $\operatorname{Spec}(A_\alpha\otimes_k B)$ pour
$A_\alpha\subset A_\beta$, ce qui résulte aussitôt de ce que $B$ est un
$k$-module plat, de (2.3.7, (ii)) et du fait que $A_\alpha$ et $A_\beta$
sont des anneaux intègres.
\end{proof}

\begin{proposition}[6.14.4]
\EGAReaderTarget{ega4.c04.prop.6-14-4}{IV.2.s6.14.4-fr}
Soient $A$ un anneau noethérien, $A'$ une $A$-algèbre noethérienne telle
que le morphisme $\operatorname{Spec}(A')\to\operatorname{Spec}(A)$ soit
normal. Soient $B$ une $A$-algèbre, $C$ la fermeture intégrale de $A$ dans
$B$; on pose $B'=B\otimes_A A'$, $C'=C\otimes_A A'$, $C'$ s'identifiant à
un sous-anneau de $B'$; alors $C'$ est la fermeture intégrale de $A'$ dans
$B'$.
\end{proposition}

\begin{proof}
Nous procéderons en plusieurs étapes.

I) \emph{Réduction au cas où l'anneau $B$ est réduit.} --- Posons
$B_0=B_{\mathrm{red}}=B/\mathfrak N$, où $\mathfrak N$ est le nilradical
de $B$, et soit $C_0$ la fermeture intégrale de $A$ dans $B_0$. On a le
lemme suivant:

\begin{lemma}[6.14.4.1]
\label{IV.2.s6.14.4.1-fr}
Soient $A$ un anneau, $B$ une $A$-algèbre, $B_0=B/\mathfrak n$ le quotient
de $B$ par un nilidéal. Si $C_0$ est la fermeture intégrale de $A$ dans
$B_0$, l'image réciproque $C$ de $C_0$ par l'homomorphisme canonique
$\varphi:B\to B_0$ est la fermeture intégrale de $A$ dans $B$.
\end{lemma}

\begin{proof}
En effet, si $x\in B$ est tel que $\varphi(x)$ vérifie une équation de
dépendance intégrale à coefficients dans $A$, on en déduit que $x$ vérifie
une relation de la forme
$x^n+a_1x^{n-1}+\cdots+a_n\in\mathfrak n$ avec $a_k\in A$, d'où, en
élevant à une puissance assez grande, une équation de dépendance intégrale
pour $x$, à coefficients dans $A$.
\end{proof}

On a donc la suite exacte de $A$-modules
\[
0\longrightarrow C\longrightarrow B\longrightarrow B_0/C_0
\]
d'où, par platitude, la suite exacte
\[
0\longrightarrow C'\longrightarrow B'\longrightarrow B'_0/C'_0
\]
où l'on a posé $B'_0=B_0\otimes_A A'$, $C'_0=C_0\otimes_A A'$. Si l'on
prouve que $C'_0$ est la fermeture intégrale de $A'$ dans $B'_0$, le lemme
(6.14.4.1) prouvera que $C'$ est la fermeture intégrale de $A'$ dans $B'$.

II) \emph{Réduction au cas où $B$ est une $A$-algèbre intègre de type fini
contenant $A$.} --- Soit $(B_\alpha)$ la famille filtrante croissante des
sous-$A$-algèbres de type fini de $B$, et soit $C_\alpha$ la fermeture
\oldpage[IV-2]{175}
intégrale de $A$ dans $B_\alpha$. Il résulte aussitôt de la définition de
la fermeture intégrale que $C$ est réunion des $C_\alpha$; si l'on pose
$B'_\alpha=B_\alpha\otimes_A A'$, $C'_\alpha=C_\alpha\otimes_A A'$,
$C'_\alpha$ est contenu dans $B'_\alpha$ par platitude, et pour la même
raison $B'$ est réunion de la famille filtrante croissante des
$B'_\alpha$ et $C'$ réunion de la famille filtrante croissante des
$C'_\alpha$. Si l'on prouve que $C'_\alpha$ est la fermeture intégrale de
$A'$ dans $B'_\alpha$ il en résultera aussitôt que $C'$ est la fermeture
intégrale de $A'$ dans $B'$. On peut donc se borner au cas où $B$ est une
$A$-algèbre de type fini, donc noethérienne; soient $\mathfrak q_i$
$(1\leqslant i\leqslant n)$ ses idéaux premiers minimaux; comme $B$ est
supposé \emph{réduit}, il s'identifie à un \emph{sous-anneau} du produit $B_0$ des
$B/\mathfrak q_i$; si $C_0$ est la fermeture intégrale de $A$ dans $B_0$,
on a $C=B\cap C_0$; si l'on pose $B'_0=B_0\otimes_A A'$,
$C'_0=C_0\otimes_A A'$, on a, par platitude,
$C'=B'\cap C'_0$ ($\mathbf{0}_{\mathrm I}$, 6.1.3); il suffit donc de prouver que
$C'_0$ est la fermeture intégrale de $A'$ dans $B'_0$. Mais $C_0$ est
composé direct des $C_i$, fermetures intégrales de $A$ dans les
$B_i=B/\mathfrak q_i$; par suite $C'_0$ est composé direct des
$C'_i=C_i\otimes_A A'$ et il suffit de montrer que $C'_i$ est la fermeture
intégrale de $A'$ dans $B'_i=B_i\otimes_A A'$. On est ainsi ramené au cas
où $B$ est intègre et une $A$-algèbre de type fini; si $\mathfrak p$ est le
noyau de l'homomorphisme $A\to B$, on a aussi
$B'=B\otimes_{A/\mathfrak p}(A'/\mathfrak pA')$; comme le morphisme
$\operatorname{Spec}(A'/\mathfrak pA')\to
\operatorname{Spec}(A/\mathfrak p)$ est normal, on peut remplacer $A$ par
$A/\mathfrak p$ et $A'$ par $A'/\mathfrak pA'$, et supposer par suite que
$A\subset B$.

III) \emph{Cas où $A$ est un corps, $B$ un corps, extension de type fini de
$A$ et tel que $A$ soit algébriquement fermé dans $B$.} --- On a alors
$C=A$, et $A'$ est une $A$-algèbre noethérienne géométriquement normale,
donc composée directe d'anneaux intègres et intégralement clos $A'_i$, les
corps des fractions $K'_i$ des $A'_i$ étant des extensions \emph{séparables} de
$A$ (4.6.1). La $A'$-algèbre $B'$ est donc composée directe des
$B\otimes_A A'_i=B'_i$ et l'on peut donc se borner au cas où $A'$ est
intègre, son corps des fractions $K'$ étant une extension séparable de $A$.
Or, comme $B$ est un $A$-module plat, $B\otimes_A A'=B'$ s'identifie à un
sous-anneau de $B\otimes_A K'$; comme $K'$ est extension séparable de $A$,
$B\otimes_A K'$ est réduit (4.3.7), et comme de plus $B$ est une extension
primaire de $A$, $B\otimes_A K'$ est intègre (4.3.2); il en est par suite
de même de $B'$. Soit $L'$ le corps des fractions de $B\otimes_A K'$ (qui
est aussi celui de $B'$); c'est évidemment un corps composé $B(K')$ de ses
sous-corps $B$ et $K'$; comme $A$ est algébriquement fermé dans $B$, que
$K'$ est une extension séparable de $A$, et que $B$ et $K'$ sont
linéairement disjoints sur $A$, $K'$ est algébriquement fermé dans
$B(K')=L'$ (Bourbaki, \emph{Alg.}, chap. V, § 9, exerc. 2); \emph{a fortiori} $A'$, étant
intégralement clos, est intégralement fermé dans $L'$, donc dans $B'$, ce
qui prouve la proposition dans ce cas.

IV) \emph{Cas où $A$ est intègre, $B$ un corps, extension finie du corps
des fractions $K$ de $A$.} --- Alors
$B'=(A'\otimes_A K)\otimes_K B$ est une $B$-algèbre noethérienne
géométriquement normale, et par suite $B'$ est composé direct d'anneaux
intègres; l'anneau total des fractions $L'$ de $B'$ est alors composé
direct d'un nombre fini de corps. Cela étant, comme $B$ est le corps des
fractions de $C$, $B'$ est un anneau de fractions de $C'$ et $L'$ est donc
aussi l'anneau total des fractions de $C'$. Or (6.14.1) comme $C$ est un
anneau normal et $\operatorname{Spec}(A')\to\operatorname{Spec}(A)$ un
morphisme normal d'anneaux noethériens, $C'$ est un anneau normal; mais il
résulte de (6.14.1.1) que $C'$ est alors intégralement fermé dans $L'$, et
\emph{a fortiori} dans $B'$, et par suite est la fermeture intégrale de $A'$ dans
$B'$.

\oldpage[IV-2]{176}
V) \emph{Fin de la démonstration.} --- D'après II), on peut supposer que
$B$ est une $A$-algèbre intègre de type fini contenant $A$. Soient $K$ le
corps des fractions de $A$, $L$ celui de $B$, qui est une extension de type
fini de $K$. Soit $M$ la fermeture algébrique de $K$ dans $L$, qui est une
extension algébrique \emph{finie} de $K$; soit $C_0$ la fermeture intégrale de $A$
dans $M$, qui est aussi la fermeture intégrale de $A$ dans $L$; on a donc
$C=B\cap C_0$; si l'on pose $C'_0=C\otimes_A A'$, on a par suite
$C'=B'\cap C'_0$ par platitude ($\mathbf{0}_{\mathrm I}$, 6.1.3). Or, il résulte de
IV) que $C'_0$ est la fermeture intégrale de $A'$ dans
$M'=M\otimes_A A'$; en outre, $M'$ est un anneau noethérien et le morphisme
$\operatorname{Spec}(M')\to\operatorname{Spec}(M)$ est normal (comme on
l'a vu dans IV)); appliquant III) à $M$ et $M'$ au lieu de $A$ et $A'$ et
à $L$ au lieu de $B$, on en déduit que $M'$ est intégralement fermé dans
$L'=L\otimes_A A'$; donc $C'_0$ est la fermeture intégrale de $A'$ dans
$L'$, et $C'=C'_0\cap B'$ est la fermeture intégrale de $A'$ dans $B'$.
C.Q.F.D.
\end{proof}

\begin{corollary}[6.14.5]
\label{IV.2.s6.14.5-fr}
Soient $A$ un anneau noethérien, $A'$ une $A$-algèbre noethérienne telle
que le morphisme $\operatorname{Spec}(A')\to\operatorname{Spec}(A)$ soit
normal. Soient $B$, $C$ deux $A$-algèbres,
$\varphi:B\to C$ un $A$-homomorphisme faisant de $C$ une $B$-algèbre, $D$
la fermeture intégrale de $B$ dans $C$. Si l'on pose
$B'=B\otimes_A A'$, $C'=C\otimes_A A'$, $D'=D\otimes_A A'$, $D'$ est la
fermeture intégrale de $B'$ dans $C'$.
\end{corollary}

\begin{proof}
Soit $(B_\lambda)$ la famille filtrante croissante des sous-$A$-algèbres
de type fini de $B$, et posons $B'_\lambda=B'_\lambda\otimes_A A'$, de
sorte que $B'$ est réunion de la famille filtrante croissante des
sous-$A'$-algèbres de type fini $B'_\lambda$. Soit $D_\lambda$ la
fermeture intégrale de $B_\lambda$ dans $C$, et posons
$D'_\lambda=D_\lambda\otimes_A A'$, de sorte que $D$ est réunion des
$D_\lambda$ et $D'$ réunion des $D'_\lambda$; si nous prouvons que
$D'_\lambda$ est la fermeture intégrale de $B'_\lambda$ dans $C'$, il en
résultera que $D'$ est la fermeture intégrale de $B'$ dans $C'$. Or
$B_\lambda$ et $B'_\lambda$ sont noethériens et le morphisme
$\operatorname{Spec}(B'_\lambda)\to\operatorname{Spec}(B_\lambda)$ est
normal (6.8.2); on peut donc appliquer (6.14.4) en remplaçant $A$, $A'$ et
$B$ par $B_\lambda$, $B'_\lambda$ et $C$ respectivement, d'où le
corollaire.
\end{proof}

On peut par exemple appliquer (6.14.5) lorsque $A$ est un anneau \emph{local
excellent} et $A'$ son complété $\widehat A$, puisque dans ce cas
$\operatorname{Spec}(A')\to\operatorname{Spec}(A)$ est un morphisme
régulier (7.8.2).

\subsection{Préschémas géométriquement unibranches.}
\label{IV.2.s6.15-fr}

\begin{env}[6.15.1]
\label{IV.2.s6.15.1-fr}
Nous dirons qu'un préschéma $X$ est \emph{unibranche} (resp.
\emph{géométriquement unibranche}) en un point $x$, ou que le point $x$
est \emph{unibranche} (resp. \emph{géométriquement unibranche}) dans $X$ si l'anneau
local $\mathcal O_{X,x}$ est \emph{unibranche} (resp. \emph{géométriquement
unibranche}) (\textbf{0}, 23.2.1); nous dirons que $X$ est \emph{unibranche} (resp.
\emph{géométriquement unibranche}) s'il l'est en tout point. Il suit de cette
définition que, pour que $X$ soit unibranche (resp. géométriquement
unibranche) en un point, il faut et il suffit que $X_{\mathrm{red}}$ le soit
en ce point.
\end{env}

\begin{env}[6.15.2]
\EGAReaderTarget{ega4.c04.statement.6-15-2}{IV.2.s6.15.2-fr}
Il faut prendre garde qu'avec les définitions de (6.15.1), un \emph{anneau local}
$A$ peut être unibranche (resp. géométriquement unibranche) \emph{sans que}
$\operatorname{Spec}(A)$ \emph{le soit}; en d'autres termes, il peut se faire qu'il
y ait des idéaux premiers $\mathfrak p$ de $A$ tels que $A_\mathfrak p$ ne
soit pas unibranche; il revient aussi au même de dire que sur un préschéma
la notion de point géométriquement unibranche \emph{n'est pas stable par
générisation}. Nous allons le voir sur l'exemple suivant.
\end{env}

\oldpage[IV-2]{177}
Soient $K$ un corps algébriquement clos de caractéristique 0, $B$ l'anneau
intègre $K[U,V,W]/(U^2(U\setmin W)\setmin V^2(U+W))$
($U$, $V$, $W$ indéterminées) de sorte que $Y=\operatorname{Spec}(B)$ est
un « cône de sommet l'origine, ayant une génératrice double ». Nous
désignerons par $u$, $v$, $w$ les images de $U$, $V$, $W$ dans $B$. Soit
$R$ le corps des fractions de $B$, et considérons dans $R$ l'élément
$t=v(u+w)/u$, qui n'appartient pas à $B$; montrons que $C=B[t]$ est la
\emph{clôture intégrale} de $B$. On a en effet $t^2=u^2\setmin w^2$, donc $t$ est entier
sur $B$, et $v=tu/(u+w)$; l'anneau $C_1=K[t,u,w]$ est intégralement clos,
car il est isomorphe à
$K[T,U,W]/(T^2\setmin U^2+W^2)$ et est donc la fermeture intégrale de l'anneau
intégralement clos $K[U,W]$ dans l'extension quadratique
$K(U,W)(\sqrt{U^2+W^2})$ de son corps des fractions (Bourbaki, \emph{Alg. comm.},
chap. V, § 1, n\textsuperscript{o} 6, prop. 18). L'anneau de fractions
$K[t,u,w,1/(u+w)]$ de $C_1$ est donc aussi intégralement clos. De la même
manière, on voit que l'anneau $C_2=K[t,v,w]$ est intégralement clos, car
$t$ est racine de l'équation
$t(t\setmin v)\setmin w^2(2v\setmin t)=0$, et par suite
$K[t,v,w,1/(t\setmin v)]=K[t,\frac{tu}{u+w},w,\frac{u+w}{tw}]$ est intégralement clos. Finalement, compte tenu de ce
que $u$ et $w$ sont algébriquement indépendants sur $K$, on prouve aisément
que $C=K[t,u,v,w]=K[t,u,w,1/(u+w)]\cap K[t,v,w,1/(t\setmin v)]$,
ce qui achève de montrer que $C$ est la clôture intégrale de $B$. Il est
immédiat que si $\mathfrak m_0$ est l'idéal maximal de $B$ engendré par
$u$, $v$, $w$ (« sommet du cône ») il existe un seul idéal maximal
$\mathfrak n_0$ de $C$ au-dessus de $\mathfrak m_0$, savoir l'idéal
engendré par $t$, $u$, $v$, $w$. Si l'on pose $A=B_{\mathfrak m_0}$, on en
déduit aisément que $A'=C_{\mathfrak n_0}$ est la clôture intégrale de
$A$, qui est donc unibranche, et par suite aussi géométriquement unibranche
puisque son corps résiduel $K$ est algébriquement clos. Mais dans $A$
l'idéal premier $\mathfrak p$ engendré par $u$ et $v$ est tel que la
clôture intégrale $A_\mathfrak p[t]$ de $A_\mathfrak p$ ne soit pas un
anneau local.

Nous verrons toutefois plus loin (9.7.10) que lorsque $X$ est un préschéma
localement noethérien tel que, si $X'$ est le normalisé de
$X_{\mathrm{red}}$, le morphisme canonique $X'\to X$ soit \emph{fini} (ce qui sera
le cas si $X$ est tel que les anneaux de ses ouverts affines soient
\emph{universellement japonais} (\textbf{0}, 23.1.1)), alors l'ensemble des points $x\in X$
où $X$ est géométriquement unibranche est \emph{localement constructible}.

\begin{env}[6.15.3]
\label{IV.2.s6.15.3-fr}
Nous dirons pour abréger qu'un morphisme $f:X\to Y$ est \emph{radiciel en
un point} $y\in Y$, si $f^{-1}(y)$ est vide ou réduit à un seul point $x$
et si $\boldsymbol{k}(x)$ est une extension radicielle de $\boldsymbol{k}(y)$. Il revient donc au
même de dire que $f$ est radiciel en tous les points de $Y$ ou que $f$ est
radiciel (\textbf{I}, 3.5.8). Si
$g:f^{-1}(y)\to\operatorname{Spec}(\boldsymbol{k}(y))$ est le morphisme image réciproque
de $f$ par le changement de base
$\operatorname{Spec}(\boldsymbol{k}(y))\to Y$, il revient au même de dire que $f$ est
radiciel au point $y\in Y$, ou que $g$ est \emph{radiciel}.
\end{env}

\begin{lemma}[6.15.3.1]
\label{IV.2.s6.15.3.1-fr}
\begin{enumerate}
\item[\textup{(i)}] Soient $f:X\to Y$, $g:Y\to Z$ deux morphismes. Si $f$ est
radiciel en un point $y$ et $g$ radiciel au point $z=g(y)$, alors $g\circ f$
est radiciel au point $g(y)$. La réciproque est vraie si $f$ est surjectif.
\item[\textup{(ii)}] Soient $f:X\to Y$, $h:Y'\to Y$ deux morphismes, et soit
$f'=f_{(Y')}:X\times_Y Y'\to Y'$. Soient $y'\in Y'$ et $y=h(y')$; alors
pour que $f$ soit radiciel au point $y$, il faut et il suffit que $f'$ soit
radiciel au point $y'$.
\end{enumerate}
\end{lemma}

\begin{proof}
\oldpage[IV-2]{178}
(i) Si $f$ est radiciel au point $y$ et $g$ radiciel au point $z$,
$g^{-1}(z)$ est réduit à $y$ et
$f^{-1}(g^{-1}(z))=f^{-1}(y)$ vide ou réduit à un seul point $x$; en
outre $\boldsymbol{k}(x)$ est extension radicielle de $\boldsymbol{k}(y)$ et $\boldsymbol{k}(y)$ extension
radicielle de $\boldsymbol{k}(z)$, donc $\boldsymbol{k}(x)$ est extension radicielle de $\boldsymbol{k}(z)$.
Inversement, supposons $f$ surjectif; si $g\circ f$ est radiciel au point
$z=g(y)$, $g^{-1}(z)$ se réduit au seul point $y$, sans quoi
$f^{-1}(g^{-1}(z))$ aurait au moins deux points distincts; en outre,
$f^{-1}(y)=f^{-1}(g^{-1}(z))$ se réduit à un seul point $x$, et par
hypothèse on a $\boldsymbol{k}(z)\subset \boldsymbol{k}(y)\subset \boldsymbol{k}(x)$, et $\boldsymbol{k}(x)$ est radiciel sur
$\boldsymbol{k}(z)$, donc $\boldsymbol{k}(y)$ est radiciel sur $\boldsymbol{k}(z)$ et $\boldsymbol{k}(x)$ radiciel sur $\boldsymbol{k}(y)$.

(ii) Soient $g:f^{-1}(y)\to\operatorname{Spec}(\boldsymbol{k}(y))$,
$g':f'^{-1}(y')\to\operatorname{Spec}(\boldsymbol{k}(y'))$ les morphismes images
réciproques de $f$ et $f'$ respectivement; comme
$f'^{-1}(y')=f^{-1}(y)\otimes_{\mathbf{k}(y)}\boldsymbol{k}(y')$ (\textbf{I}, 3.6.4), $g'$ est image
réciproque de $g$ et l'assertion se réduit à (2.6.1, (v)).
\end{proof}

Soit alors $X$ un préschéma \emph{réduit} n'ayant qu'un nombre \emph{fini} de
composantes irréductibles, et soit $X'$ son \emph{normalisé} (\textbf{II}, 6.3.8); on sait
\emph{(loc. cit.)} que le morphisme canonique $f:X'\to X$ est surjectif. La
définition (6.15.1) montre alors que pour que $X$ soit géométriquement
unibranche en un point $x$, il faut et il suffit que $f$ soit \emph{radiciel en ce
point}; pour que $X$ soit géométriquement unibranche, il faut et il suffit
donc que $f$ soit \emph{radiciel}.

\begin{env}[6.15.4]
\EGAReaderTarget{ega4.c04.statement.6-15-4}{IV.2.s6.15.4-fr}
Généralisant la définition donnée dans (\textbf{I}, 2.2.9), nous dirons que, pour
deux préschémas \emph{réduits} $X$, $Y$, un morphisme $f:X\to Y$ est
\emph{birationnel} si la restriction de $f$ à l'ensemble des points
maximaux de $X$ est une \emph{bijection} de cet ensemble sur l'ensemble des points
maximaux de $Y$, et si, pour tout point maximal $y$ de $Y$, le morphisme
$X\times_Y\operatorname{Spec}(\mathcal O_y)\to\operatorname{Spec}(\mathcal O_y)$ déduit de
$f$ est un isomorphisme (autrement dit, si la fibre $f^{-1}(y)$ est réduite
à un seul point $x$ (maximal dans $X$) et si l'homomorphisme
$\boldsymbol{k}(y)\to \boldsymbol{k}(x)$ déduit de $f$ est bijectif); cette notion se réduit à celle
de (\textbf{I}, 2.2.9) lorsque $X$ et $Y$ n'ont qu'un nombre \emph{fini} de composantes
irréductibles.
\end{env}

\begin{lemma}[6.15.4.1]
\label{IV.2.s6.15.4.1-fr}
Soit $f:X\to Y$ un morphisme, $g:Y'\to Y$ un morphisme plat,
$X'=X\times_Y Y'$, $f'=f_{(Y')}:X'\to Y'$. Si $f_{\mathrm{red}}$ est
birationnel, il en est de même de $f'_{\mathrm{red}}$.
\end{lemma}

\begin{proof}
En effet, si $y'$ est un point maximal de $Y'$, on sait (2.3.4, (ii)) que
$y=g(y')$ est un point maximal de $Y$; il existe un seul point $x$ de $X$
dans $f^{-1}(y)$ par hypothèse; de plus $x$ est maximal dans $X$ et tel que
$\boldsymbol{k}(x)=\boldsymbol{k}(y)$. Il résulte donc de (\textbf{I}, 3.4.9) que $f'^{-1}(y')$ est réduit à
un seul point $x'$ et que $\boldsymbol{k}(x')=\boldsymbol{k}(y')$. On conclut le raisonnement en
remarquant que d'après (2.3.7, (ii)), toute composante irréductible de $X'$
domine une composante irréductible de $Y'$.
\end{proof}

\begin{proposition}[6.15.5]
\label{IV.2.s6.15.5-fr}
Soit $f:X\to Y$ un morphisme tel que $f_{\mathrm{red}}$ soit entier et
birationnel \textup{(6.15.4)}, donc surjectif.
\begin{enumerate}
\item[\textup{(i)}] Pour que $Y$ soit géométriquement unibranche en un point $y$,
il faut et il suffit que $f$ soit radiciel au point $y$ et que $X$ soit
géométriquement unibranche en l'unique point $x$ de $f^{-1}(y)$.
\item[\textup{(ii)}] Pour que $Y$ soit géométriquement unibranche, il faut et il
suffit que $X$ soit géométriquement unibranche et que $f$ soit radiciel.
\end{enumerate}
\end{proposition}

\begin{proof}
On peut évidemment se borner au cas où $X$ et $Y$ sont \emph{réduits}; le fait
que $f$ soit \emph{surjectif} résulte de ce que $f$ est un morphisme fermé (\textbf{II},
6.1.10) et que $f(X)$ contient par hypothèse tous les points maximaux de
$Y$. Cela montre aussitôt que (i) entraîne (ii). Utilisant (\textbf{I}, 3.6.5) et
(\textbf{II}, 6.1.5), on peut supposer, pour prouver (i), que
$Y=\operatorname{Spec}(\mathcal O_{Y,y})$,
\oldpage[IV-2]{179}
autrement dit
$Y=\operatorname{Spec}(A)$, où $A$ est \emph{local}. Comme $f$ est affine, on a
aussi $X=\operatorname{Spec}(B)$. Si $A$ est géométriquement unibranche,
il est \emph{intègre}, donc il en est de même de $B$ puisque $f$ est birationnel;
inversement, si $f$ est radiciel au point $y$ et $X$ géométriquement
unibranche au point $x$, $B$ n'a qu'un seul idéal maximal (puisqu'il est
entier sur $A$ et que tout idéal maximal de $B$ est donc au-dessus de celui
de $A$), autrement dit $B$ est un anneau \emph{local}, et dire que $X$ est
géométriquement unibranche au point $x$ signifie que $B$ est
géométriquement unibranche, donc \emph{intègre}. Comme $f$ est dominant par
hypothèse et que $A$ est réduit, on en conclut que $A\subset B$ (\textbf{I},
1.2.7), donc $A$ est aussi \emph{intègre}. Ainsi, pour prouver (i), on peut se
borner au cas où $A$ et $B$ sont locaux, intègres, $A$ étant contenu dans
$B$ et ayant même corps des fractions, et par suite (puisque $B$ est entier
sur $A$) même \emph{clôture intégrale} $C$. Mais alors l'assertion résulte de
(6.15.3.1, (i)) appliqué aux morphismes
$\operatorname{Spec}(C)\to\operatorname{Spec}(B)$ et
$\operatorname{Spec}(B)\to\operatorname{Spec}(A)$, vu l'interprétation
donnée dans (6.15.3) de la propriété d'être géométriquement unibranche en
un point.
\end{proof}

\begin{proposition}[6.15.6]
\label{IV.2.s6.15.6-fr}
Soient $k$ un corps, $X$ un $k$-préschéma. Si $X$ est normal, alors, pour
toute extension $k'$ de $k$, $X'=X\otimes_k k'$ est géométriquement
unibranche.
\end{proposition}

\begin{proof}
On sait que $k'$ est extension algébrique d'une extension pure $k_0$ de
$k$, et si $k''$ est la plus grande extension séparable de $k_0$ contenue
dans $k'$, $k'$ est une extension radicielle de $k''$ et $k''$ extension
séparable de $k$. On sait (6.14.2) que $X''=X\otimes_k k''$ est normal;
comme $X\otimes_k k'=X''\otimes_{k''}k'$, on voit qu'on peut se borner au
cas où l'extension $k'$ de $k$ est \emph{radicielle}. En outre (\textbf{I}, 3.6.5), on peut
supposer que $X=\operatorname{Spec}(A)$, où $A$ est un anneau \emph{local intègre
et intégralement clos} (puisque $X$ est normal). Le morphisme projection
$f:X'\to X$ est un homéomorphisme, puisque
$\operatorname{Spec}(k')\to\operatorname{Spec}(k)$ est un homéomorphisme
universel (2.4.5); comme $X'=\operatorname{Spec}(A')$ où
$A'=A\otimes_k k'$, on voit donc que $A'$ est un anneau local dont le
nilradical $\mathfrak N'$ est l'unique idéal premier minimal, d'où
$X'_{\mathrm{red}}=\operatorname{Spec}(A'_0)$, où
$A'_0=A'/\mathfrak N'$ est un anneau local intègre; en outre, si $K$ est
le corps des fractions de $A$, le corps des fractions $K'_0$ de $A'_0$ est
radiciel sur $K$, puisque le morphisme $f$ est radiciel. Comme $A'_0$ est
entier sur $A$, sa clôture intégrale $B$ est aussi la fermeture intégrale
de $A$ dans $K'_0$. Mais comme $A$ est intégralement clos, on sait
(Bourbaki, \emph{Alg. comm.}, chap. V, § 2, n\textsuperscript{o} 3, lemme 3) que $B$ est
l'ensemble des $x\in K'_0$ dont une puissance $p^m$-ème (pour un $m$ assez
grand) appartient à $A$ ($p$ étant l'exposant caractéristique de $K$);
en outre il n'existe qu'un seul idéal premier de $B$ au-dessus de tout idéal
premier de $A$; en particulier $B$ est un anneau local et son corps résiduel
est une extension radicielle de celui de $A$, et \emph{a fortiori} de celui de
$A'_0$, ce qui prouve que $A'_0$ est géométriquement unibranche, et par
suite il en est de même de $X'$.
\end{proof}

\begin{proposition}[6.15.7]
\EGAReaderTarget{ega4.c04.prop.6-15-7}{IV.2.s6.15.7-fr}
Soient $k$ un corps, $X$ un $k$-préschéma, $k'$ une extension de $k$,
$X'=X\otimes_A k'$. Soient $x'$ un point de $X'$, $x$ sa projection dans
$X$. Pour que $X$ soit géométriquement unibranche au point $x$, il faut et
il suffit que $X'$ soit géométriquement unibranche au point $x'$. Pour que
$X$ soit géométriquement unibranche, il faut et il suffit que $X'$ le soit.
\end{proposition}

\begin{proof}
La seconde assertion résulte de la première et du fait que la projection
$f:X'\to X$ est un morphisme surjectif. Pour démontrer la première
assertion, on peut, en vertu de (\textbf{I}, 5.1.8), se borner au cas où $X$ est
réduit, et en vertu de (\textbf{I}, 3.6.5), au cas
\oldpage[IV-2]{180}
où $X=\operatorname{Spec}(A)$ (avec $A=\mathcal O_{X,x}$) est un schéma local.
Notons que $A'=\mathcal O_{X',x'}$ est par hypothèse un $A$-module
fidèlement plat, donc contenant $A$, et par suite l'hypothèse que $X$ est
géométriquement unibranche au point $x$, ou l'hypothèse que $X'$ est
géométriquement unibranche au point $x'$, entraînent toutes deux que $A$
est un anneau local intègre (puisque $A$, étant réduit, est aussi isomorphe
à un sous-anneau de $A'_{\mathrm{red}}$). Soient $K$ le corps des fractions
de $A$, $B$ la clôture intégrale de $A$, et posons
$Y=\operatorname{Spec}(B)$; il revient au même de dire que $X$ est
géométriquement unibranche au point $x$ ou que le morphisme $g:Y\to X$
est radiciel au point $x$ (6.15.3). Soit
$Y'=Y\otimes_k k'=Y\times_X X'$ de sorte que l'on a le diagramme
commutatif
\[
\xymatrix{
Y \ar[d]_g & Y' \ar[l]^h \ar[d]_{g'}\\
X & X' \ar[l]_f
}
\]
Notons que $f$ est un morphisme plat; donc (6.15.4.1) le morphisme entier
$g'_{\mathrm{red}}$ est birationnel. D'autre part, comme $Y$ est normal,
$Y'$ est géométriquement unibranche (6.15.6). Pour que $X'$ soit
géométriquement unibranche en $x'$, il faut et il suffit donc que $g'$
soit radiciel au point $x'$ (6.15.5). Or, $x$ est la projection de $x'$ dans
$X$ et il est équivalent de dire que $g$ est radiciel au point $x$ ou que
$g'$ est radiciel au point $x'$ (6.15.3.1); enfin, pour que $X$ soit
géométriquement unibranche au point $x$, il faut et il suffit que $g$ soit
radiciel en ce point, d'où la proposition.
\end{proof}

\begin{lemma}[6.15.8]
\label{IV.2.s6.15.8-fr}
Soient $k$ un corps séparablement clos (autrement dit, tel que la clôture
algébrique de $k$ soit radicielle sur $k$), $X$ un $k$-préschéma localement
de type fini sur $k$, $x$ un point fermé de $X$. Si $X$ est unibranche au
point $x$, il est géométriquement unibranche en ce point.
\end{lemma}

\begin{proof}
En effet, on sait (I, 6.4.2) que $k(x)$ est une extension algébrique de
$k$; comme le corps résiduel de la clôture intégrale de
$(\mathcal O_{X,x})_{\mathrm{red}}$ (clôture intégrale qui est par
hypothèse un anneau local) est une extension algébrique de $k(x)$, c'est
une extension radicielle de $k$ par hypothèse, donc aussi de $k(x)$.
\end{proof}

\begin{corollary}[6.15.9]
\label{IV.2.s6.15.9-fr}
Soient $k$ un corps, $X$ un $k$-préschéma localement de type fini sur $k$,
$k'$ une extension séparablement close de $k$ (autrement dit, telle que la
clôture algébrique de $k'$ soit radicielle sur $k'$); alors, pour que $X$
soit géométriquement unibranche, il faut et il suffit que
$X'=X\otimes_k k'$ soit unibranche.
\end{corollary}

\begin{proof}
Compte tenu de (6.15.7), on est ramené à prouver que si $X$ est unibranche
et $k$ séparablement clos, alors $X$ est géométriquement unibranche. Or, $X$
est géométriquement unibranche en ses points fermés (6.15.8); nous en
conclurons dans (10.4.9) que $X$ est géométriquement unibranche en tous ses
points, en nous appuyant sur le fait (signalé à la fin de (6.15.2)) que
l'ensemble des points où $X$ est géométriquement unibranche est
constructible (bien entendu, le cor. (6.15.9) ne sera pas utilisé
jusque-là).
\end{proof}

Ce résultat justifie, en une certaine mesure, la terminologie de
« géométriquement unibranche ».

